% --- Farbdefinitionen ---
\definecolor{codebg}{RGB}{247,247,247}
\definecolor{terminalbg}{RGB}{235, 240, 245}
\definecolor{terminalfg}{RGB}{45, 45, 45}
\definecolor{commentgreen}{RGB}{0,128,0}
\definecolor{keywordblue}{RGB}{0,0,255}
\definecolor{stringstring}{RGB}{163,21,21}

% --- Kopf- und Fußzeilen (fancyhdr) ---
\pagestyle{fancy}
\fancyhf{}
\renewcommand{\sectionmark}[1]{\markboth{#1}{}}
\lhead{\leftmark}
\rhead{\MyTitle} % Nutzt jetzt den automatischen Titel des Pakets
\rfoot{Seite \thepage\ von \pageref{LastPage}}
\renewcommand{\headrulewidth}{0.4pt}
\renewcommand{\footrulewidth}{0.4pt}

% --- siunitx Setup ---
\sisetup{locale = DE, per-mode = symbol}

% --- Listings (Code Styles) ---
% Standard-Style für echten Quellcode
\lstset{
    backgroundcolor=\color{codebg},
    commentstyle=\color{commentgreen},
    keywordstyle=\color{keywordblue}\bfseries,
    numberstyle=\tiny\color{gray},
    stringstyle=\color{stringstring},
    basicstyle=\ttfamily\small,
    breakatwhitespace=false,
    breaklines=true,
    captionpos=b,
    keepspaces=true,
    numbers=left,
    numbersep=8pt,
    showspaces=false,
    showstringspaces=false,
    showtabs=false,
    tabsize=4,
    frame=single,
    rulecolor=\color{lightgray}
}

% Style für Terminal-Ausgaben
\lstdefinestyle{terminal}{
    backgroundcolor=\color{terminalbg},
    basicstyle=\ttfamily\small\color{terminalfg},
    numbers=none,            
    frame=single,
    rulecolor=\color{gray},
    breaklines=true,
    keepspaces=true,
    commentstyle=\color{terminalfg},
    keywordstyle=\color{terminalfg},
    stringstyle=\color{terminalfg}
}

% Eigener Kurzbefehl für Inline-Code
\newcommand{\code}[1]{\colorbox{codebg}{\ttfamily\small #1}}
